% IEEE journal template for texforge. Text written in double braces is a
% placeholder that is filled in when the project is created; the % comments
% say what to replace.
\documentclass[10pt,journal]{IEEEtran}
\usepackage[{{language}}]{babel}
\addto\captionsspanish{\renewcommand{\tablename}{Tabla}}
\usepackage[T1]{fontenc}
\usepackage{amsmath,amssymb}
\usepackage{booktabs}
\usepackage{float}
\usepackage{graphicx}
\usepackage{url}
\usepackage[pdfauthor={{{author}}},pdftitle={{{title}}},colorlinks,linkcolor=blue,citecolor=blue,urlcolor=blue]{hyperref}

\begin{document}

  % Title of the article; keep it short enough to fit two lines.
  \title{{{title}}}

  % Name, italic affiliation and monospaced e-mail of the author.
  \author{
    {{{author}}}\\
    \textit{{{affiliation}}}\\
    \texttt{{{email}}}
  }

  \maketitle
  % Running head of the inner pages, taken from the {{venue}} placeholder.
  \markboth{{{venue}}}{}

  \begin{abstract}
    % Replace this text with a short summary of the problem, the method, the
    % main result and why the result matters. Three or five sentences are enough.
    This template shows the shape of an IEEE journal paper. Replace this text
    with a short summary of the problem you address, the method you use, the
    main result you obtain, and why the result matters. A reader should grasp
    the contribution without having to read the whole document.
  \end{abstract}

  \begin{IEEEkeywords}
    % Comma separated index terms; they come from the {{keywords}} placeholder.
    {{keywords}}
  \end{IEEEkeywords}

  \section{Introduction}
  % State the problem, why it matters, what this document contributes, and
  % how the rest of the paper is organised.
  \IEEEPARstart{T}{his} document is a complete example of an IEEE journal
  paper. Replace this text with the context of your work: the problem you
  attack, the evidence that the problem matters, and the contribution of the
  paper. Cite the foundational work of the field here~\cite{doe2019}. Two or
  three short paragraphs are usually enough, and the last one should point the
  reader to the sections that follow.

  The sample table, the sample equation and the two sample diagrams are part
  of the template: keep them while you draft, then replace each one with your
  own content. The two column layout, the running head and the reference list
  take care of themselves once the template fields are filled in.

  \section{Related Work}
  % Summarise the closest prior work, one paragraph per line of research, and
  % name the gap that your document fills.
  Summarise here the work closest to yours and how it differs from your own
  approach~\cite{jones2023}. A textbook reference supports the standard method
  that you build on~\cite{williams2019}. End the section by naming the gap
  that the present document fills.

  Keep the list short: two or three lines of research are usually enough to
  position your own work. Every source discussed here should also be cited in
  the sections that follow, so that the reference list stays complete.

  \section{Data and Methods}
  % Describe where the data comes from, how it was prepared, and how the
  % models were fitted and evaluated. Cite the data set that you used.
  The sample data used in this template is synthetic and only shows the layout
  of this section; replace it with a description of your own data and cite the
  data set~\cite{openlab2024}. The error of a model is defined as
  \begin{equation}
    E(\theta) = \frac{1}{n} \sum_{i=1}^{n} \bigl( y_i - f_\theta(x_i) \bigr)^2
    \label{eq:error}
  \end{equation}
  where $y_i$ is the observed value and $f_\theta(x_i)$ is the prediction of
  the model for sample $i$. Equation~\eqref{eq:error} is reported for every
  variant in Table~\ref{tab:scores}.

  Write also how the data was prepared, how it was split into a training part
  and a test part, and which settings were used for each run. A reader who
  wants to repeat the experiment needs that detail, and the software used to
  run it deserves a citation of its own.

  \section{Results and Discussion}
  % Report the main numbers first, then explain what they mean. Replace the
  % table with your own measurements and keep the best value in bold.
  Table~\ref{tab:scores} compares three variants of the method on the same
  sample data. The best error comes from the full method, at the cost of a
  longer runtime; say here whether that trade off is acceptable for the
  problem you study.

  \begin{table}[H]
    \centering
    \caption{Error and runtime of each variant on the sample data}
    \label{tab:scores}
    \begin{tabular}{lcc}
      \toprule
      Variant & Error & Runtime (s) \\
      \midrule
      Baseline & 0.184 & 12.4 \\
      Improved & 0.151 & 15.9 \\
      Full method & \textbf{0.132} & 18.6 \\
      \bottomrule
    \end{tabular}
  \end{table}

  Read the table before the diagrams: which variant wins, by how much, and
  whether the difference matters for the problem at hand. Keep the discussion
  tied to the numbers instead of repeating them, and let each figure carry the
  one point that the text cannot show. The steps of the pipeline run from left
  to right in the first figure below.

  % Replace the two diagrams below with your own pipeline and your own model
  % diagram. The options are the style preset, the placement (pos=H keeps the
  % drawing where it is written), the caption and the width. Keep each drawing
  % close to square: a tall one is stretched to the width of the column and
  % ends up much larger than the text around it.
  \begin{mermaid}[style=technical, pos=H, caption={Pipeline used to obtain the results}, width=\linewidth]
    flowchart LR
    collect[Collect] --> prepare[Clean] --> train[Fit] --> report[Report]
  \end{mermaid}

  The second figure draws what the reported error depends on: the fitted model,
  the observed value of the sample and the settings of the run. Both diagrams
  are drawn from their source when the document is built, so the project never
  carries image files.

  \begin{graphviz}[style=editorial, pos=H, caption={Terms that the reported error depends on}, width=\linewidth]
    digraph G {
      node [shape=box]
      sample -> model
      settings -> model
      model -> error
      error [shape=ellipse]
    }
  \end{graphviz}

  \section{Conclusions}
  % Restate the main result in one sentence, name the open limitations and
  % the next step.
  Replace this text with a short closing: what the results show, which
  limitations remain, and what should be done next. Name the data or the
  people that made the work possible, and state clearly what the next study
  should test. A final citation to the wider context of the problem is
  welcome~\cite{openlab2024}.

  \bibliographystyle{IEEEtran}
  \bibliography{bib/references}

  \appendices
  \section{Key code}
  % Keep only the few lines that a reader needs to reproduce the result, and
  % keep them short enough to fit the width of one column.
  The listing below is the whole of the code behind Table~\ref{tab:scores}.

  \begin{code}[lang=python, numbers=true]
def error(obs, pred):
    """Mean squared error."""
    total = 0.0
    for o, p in zip(obs, pred):
        total += (o - p) ** 2
    return total / len(obs)


def report(scores):
    for name in sorted(scores):
        value = scores[name]
        print(name, round(value, 3))
  \end{code}

\end{document}
